% Part VI: Applications and Future Directions
% This file is \include'd from main.tex

\chapter{LLM-Assisted Software Synthesis}
\label{ch:llm-synthesis}

Large Language Models (LLMs) generate plausible code with remarkable
fluency, yet their outputs carry no structural guarantees: type errors
slip through, modules fail to compose, and subtle semantic bugs hide
behind syntactically correct surfaces.  This chapter shows how
Judgment Geometry supplies the missing \emph{verification and repair
layer} for LLM-generated code.

%% ─── 27.1 ──────────────────────────────────────────────────────
\section{The structural deficit of LLM code}
\label{sec:llm-deficit}

An LLM treats code generation as next-token prediction: given a prompt
$p$, it samples a continuation $c$ from a learned distribution
$P(c \mid p)$.  Nothing in this process enforces that:
\begin{enumerate}[label=(\roman*)]
  \item $c$ type-checks against the project's type system,
  \item $c$ is consistent with specifications in \emph{other} modules,
  \item $c$ satisfies non-local invariants (security policies,
        performance contracts, API versioning).
\end{enumerate}
Testing catches some of these failures, but testing is inherently
\emph{local}: a test suite checks individual behaviors ($H^0$, in our
language) without examining whether local sections \emph{glue} into a
global one ($H^1$).  Formal verification, at the other extreme, is
too rigid---it demands binary pass/fail and cannot exploit partial
correctness.

Judgment Geometry occupies the middle ground: it is \emph{graded}
(partial correctness is meaningful), \emph{cohomological} (it detects
non-local inconsistencies), and \emph{iterative} (the descent algorithm
produces a convergent repair sequence).

%% ─── 27.2 ──────────────────────────────────────────────────────
\section{JuGeo as LLM supervisor}
\label{sec:llm-supervisor}

We instantiate the four objects for LLM-assisted development.

\begin{definition}[LLM Development Site]
\label{def:llm-site}
Let $\mathcal{P}$ be a software project with modules
$M_1,\dots,M_n$.  Define the \emph{LLM development site}
$\STop_{\mathrm{LLM}}$ as follows:
\begin{itemize}
  \item \textbf{Objects}: modules, files, functions, and test suites
        of~$\mathcal{P}$, ordered by containment.
  \item \textbf{Morphisms}: inclusion and dependency maps
        (imports, function calls, type references).
  \item \textbf{Topology}: the coverage $J$ in which a family
        $\{U_i \to U\}$ covers~$U$ when every public symbol of~$U$
        is exercised by at least one~$U_i$.
\end{itemize}
\end{definition}

\begin{definition}[LLM Proof Structure]
\label{def:llm-pstr}
The presheaf $\PStr_{\mathrm{LLM}}$ assigns to each object $U$:
\begin{itemize}
  \item A \emph{specification} $A_U$ (type signature, docstring
        contract, or formal pre/post-condition).
  \item An \emph{evidence bundle} $E_U$: the LLM-generated code
        together with any passing tests.
\end{itemize}
A section $s \in \PStr_{\mathrm{LLM}}(U)$ is a pair $(A_U, E_U)$
asserting ``evidence $E_U$ satisfies specification~$A_U$\!.''
\end{definition}

The \emph{Cohomological Spectrum} $\CSpec_{\mathrm{LLM}}$ is the
\v{C}ech complex of $\PStr_{\mathrm{LLM}}$.  Concretely:
\[
  \Cech^0 = \prod_i \PStr(U_i), \qquad
  \Cech^1 = \prod_{i<j} \PStr(U_i \cap U_j),
\]
where $U_i \cap U_j$ is the shared interface between modules $M_i$
and~$M_j$ (exported types, shared state, API contracts).  A class
$[\alpha] \in \check{H}^1(\STop_{\mathrm{LLM}}, \PStr_{\mathrm{LLM}})$
with $[\alpha]\neq 0$ witnesses an \emph{integration bug}: the
LLM-generated code for $M_i$ and $M_j$ individually type-checks but
disagrees on the overlap.

\begin{example}
Module $M_1$ exports a function \texttt{get\_user} returning a
dictionary with key \texttt{"name"}.  The LLM generates $M_2$
expecting a key \texttt{"username"}.  Both modules pass their local
tests (the LLM generates a mock with the expected key), but the
overlap section records the mismatch.  This is a non-trivial class in
$\check{H}^1$.
\end{example}

The \emph{Graded Evaluation} $\GEval_{\mathrm{LLM}}$ grades each
section along multiple axes:
\begin{equation}
\label{eq:llm-grade}
  g(s) \;=\; g_{\mathrm{type}} \;\otG\;
             g_{\mathrm{test}} \;\otG\;
             g_{\mathrm{style}} \;\otG\;
             g_{\mathrm{perf}},
\end{equation}
where each component lies in $\Grade = (-\infty, 0]$ and
$\otG = +$ aggregates them (recall $\gone = 0$ is perfect).

%% ─── 27.3 ──────────────────────────────────────────────────────
\section{The descent loop for LLM code}
\label{sec:llm-descent}

The central observation is that \emph{iterative LLM repair is exactly
descent}.

\begin{construction}[LLM Descent Algorithm]
\label{con:llm-descent}
Given an LLM $\mathcal{L}$, a project site $\STop$, and a
specification presheaf $\PStr$:
\begin{enumerate}
  \item \textbf{Generate.}\; For each $U_i$, prompt $\mathcal{L}$
        to produce a candidate section
        $s_i^{(0)} \in \PStr(U_i)$.
  \item \textbf{Compute coboundary.}\; Form
        $\delta^0 s^{(0)} \in \Cech^1$; identify all non-zero
        components $\alpha_{ij}$.
  \item \textbf{Repair.}\; For each obstruction $\alpha_{ij}$,
        construct a \emph{repair prompt}
        $p_{ij} = \texttt{``Fix: module } M_i \texttt{ expects \dots\ but }
        M_j \texttt{ provides \dots''}$.
        Feed $p_{ij}$ to~$\mathcal{L}$; obtain $s_i^{(1)}$
        or~$s_j^{(1)}$.
  \item \textbf{Re-evaluate.}\; Compute $\delta^0 s^{(1)}$.
        If $\check{H}^1 = 0$ and
        $\GEval(s^{(1)}) \geq g_{\min}$, accept.
        Otherwise return to step~3.
  \item \textbf{Budget.}\; If the iteration count exceeds a bound
        $N$ or the grade fails to improve for $k$ consecutive
        iterations, halt and report the residual obstructions.
\end{enumerate}
\end{construction}

\begin{proposition}
\label{prop:llm-descent-mono}
If each repair step strictly improves the grade of at least one
obstruction component (i.e., $g(\alpha_{ij}^{(t+1)}) >
g(\alpha_{ij}^{(t)})$ for some~$(i,j)$), and the grade is bounded
above by~$\gone$, then the descent loop terminates in at most
$n(n-1)/2 \cdot \lceil (\gone - g_{\min})/\epsilon \rceil$ steps,
where $\epsilon$ is the minimum per-step improvement.
\end{proposition}

\begin{proof}
There are at most $\binom{n}{2}$ overlap pairs.  Each
can improve at most $(\gone - g_{\min})/\epsilon$ times before
reaching~$\gone$.  Once every overlap grade reaches~$\gone$, the
coboundary vanishes and $\check{H}^1 = 0$.
\end{proof}

\begin{remark}
In practice one does not require $\check{H}^1 = 0$ exactly; a graded
threshold $g_{\min}$ suffices.  This is a key advantage of graded
cohomology over binary verification: the system can declare
``good enough'' and defer minor mismatches.
\end{remark}

%% ─── 27.4 ──────────────────────────────────────────────────────
\section{Prompt engineering as site construction}
\label{sec:prompt-site}

We now observe that the art of prompt engineering admits a precise
geometric interpretation.

\begin{proposition}[Prompts as Covers]
\label{prop:prompt-cover}
Let $\mathcal{T}$ be the target specification (the ``task'' given
to the LLM).  A \emph{prompt decomposition}
$p = \{p_1, \dots, p_k\}$ induces a covering family
$\{U_1, \dots, U_k\}$ of~$\mathcal{T}$ in $\STop$, where each
$U_i$ is the sub-task addressed by prompt~$p_i$.  The LLM's response
to $p_i$ is a local section $s_i \in \PStr(U_i)$.  The full
response glues iff $\check{H}^1(\{U_i\}, \PStr) = 0$.
\end{proposition}

Under this interpretation:
\begin{itemize}
  \item \textbf{Few-shot examples} are \emph{covering sieves}:
        they demonstrate the desired behavior on specific
        sub-regions of the task, guiding the LLM toward sections
        that agree on overlaps.
  \item \textbf{The context window} is the \emph{open cover}
        $\mathcal{U}$: the finite amount of information the LLM can
        attend to at once.
  \item \textbf{Prompt decomposition} is \emph{refinement of the
        cover}: splitting a complex prompt into sub-prompts
        corresponds to refining $\mathcal{U}$ into a finer cover,
        which can only \emph{improve} the cohomological condition
        (by Leray's theorem, refinement does not increase~$H^1$).
  \item \textbf{Chain-of-thought prompting} is \emph{choosing a good
        nerve}: it makes the simplicial structure of the cover
        explicit, so the LLM can track overlaps.
\end{itemize}

%% ─── 27.5 ──────────────────────────────────────────────────────
\section{Multi-agent orchestration via the fibration}
\label{sec:multi-agent}

Modern LLM systems deploy multiple agents, each responsible for one
module or aspect of the codebase.  JuGeo provides the orchestration
framework.

\begin{construction}[Multi-Agent Fibration]
\label{con:multi-agent}
Let $\{A_1, \dots, A_n\}$ be LLM agents, with agent $A_i$
responsible for module~$M_i$.  The \emph{judgment fibration}
$\pi\colon \mathcal{E} \to \STop$ has fiber $\pi^{-1}(U_i)$
consisting of all code artifacts agent~$A_i$ can produce for~$U_i$.
The orchestrator's job is to select a \emph{global section}
$s \in \Gamma(\STop, \mathcal{E})$---i.e., a choice of artifact per
module that glues across all overlaps.  This is precisely the
descent problem.
\end{construction}

The orchestrator proceeds as follows:
\begin{enumerate}
  \item Each agent $A_i$ proposes a local section $s_i$.
  \item The orchestrator computes $\delta^0 s$ and identifies
        obstructions $\alpha_{ij}$.
  \item For each obstruction, the orchestrator sends a
        \emph{compatibility constraint} to agents $A_i$ and~$A_j$.
  \item The agents regenerate; the process iterates.
\end{enumerate}

This is the exact pattern used by systems such as ChatDev and
MetaGPT, but JuGeo gives it a \emph{mathematical foundation}:
convergence, termination, and quality guarantees follow from the
general descent theory of Part~II.

%% ─── 27.6 ──────────────────────────────────────────────────────
\section{Application: full-stack web generation}
\label{sec:webapp-gen}

We illustrate the framework on a concrete application:
generating a complete web application (HTML, CSS, JavaScript,
backend~API).

\begin{example}[Web Application as Stratified Site]
\label{ex:webapp-site}
A web application defines a natural three-stratum site, analogous to
the linguistic strata of Part~III:
\begin{center}
\begin{tabular}{lll}
\textbf{Stratum} & \textbf{Web layer} & \textbf{Linguistic analogue} \\
\hline
$\sigma_1$ & HTML (structure) & $\synt$ (syntax) \\
$\sigma_2$ & CSS (presentation) & $\phon$ (phonology / surface form) \\
$\sigma_3$ & JS (behavior) & $\sem$ (semantics) \\
$\sigma_4$ & API (data) & $\prag$ (pragmatics / context)
\end{tabular}
\end{center}
Cross-stratum descent asks: does the JavaScript match the HTML
structure?  Do the CSS selectors reference elements that actually
exist?  Does the API return data in the format the frontend expects?
\end{example}

The generation loop for a web app is:
\begin{enumerate}
  \item \textbf{Specification.}\;  The user provides a natural-language
        description $\mathcal{T}$.
  \item \textbf{Site construction.}\;  Parse $\mathcal{T}$ into a
        module decomposition $\{U_i\}$ (pages, components, API
        endpoints).
  \item \textbf{LLM generation.}\;  For each $(U_i, \sigma_j)$
        (module $\times$ stratum), the LLM generates a local section.
  \item \textbf{Cross-stratum descent.}\;  Compute
        $\check{H}^1$ across strata: verify that CSS classes match
        HTML elements, that JS event handlers reference existing DOM
        nodes, that API calls match endpoint signatures.
  \item \textbf{Repair.}\;  Feed obstructions back to the LLM as
        repair prompts.
  \item \textbf{Compilation and testing.}\;  Run the project's build
        system and test suite; feed results into $\GEval$.
  \item \textbf{Iterate.}\;  Repeat until $\GEval \geq g_{\min}$.
\end{enumerate}

\begin{remark}
The \texttt{jugeo-webapp} module in the JuGeo codebase implements
a prototype of this loop, using the descent engine from
\texttt{src/jugeo/} to orchestrate HTML/CSS/JS generation and
cross-layer verification.
\end{remark}

%% ─── 27.7 ──────────────────────────────────────────────────────
\section{Comparison with existing approaches}
\label{sec:llm-comparison}

We compare JuGeo-supervised LLM code generation with three
alternatives.

\medskip\noindent
\textbf{Plain LLM generation.}\;
Without JuGeo, an LLM has no mechanism to detect non-local
inconsistencies.  It can only ``hope'' that independently generated
modules cohere.  JuGeo adds the cohomological check ($\CSpec$)
and the descent repair loop, converting hope into a convergent
algorithm.

\medskip\noindent
\textbf{Formal verification.}\;
Systems like Dafny and Lean provide strong guarantees but require
full formal specifications and produce binary verdicts.
JuGeo is \emph{graded}: a section with
$g(s) = -0.3$ is meaningful (``almost correct''),
enabling iterative improvement.  Moreover, JuGeo can incorporate
formal verification as one component of $\GEval$ (set
$g_{\mathrm{formal}} = \gone$ if the prover succeeds,
$g_{\mathrm{formal}} = \gzero$ otherwise), subsuming
formal methods as a special case.

\medskip\noindent
\textbf{Testing alone.}\;
A test suite checks $H^0$: does each module individually behave
correctly?  JuGeo also checks $H^1$: do modules compose correctly?
This is precisely the gap that integration bugs exploit.
By computing the full \v{C}ech complex, JuGeo detects classes
of bugs that no finite test suite can reliably catch---those
arising from the \emph{interaction} of correct components.

\begin{proposition}[Subsumption]
\label{prop:subsumption}
Let $\mathcal{V}$ be any verification methodology that assigns a
binary verdict $v_U \in \{0,1\}$ to each module~$U$.  Then
$\mathcal{V}$ embeds into $\GEval$ via $g_U = v_U \cdot \gone$,
and the $\check{H}^1$ computation of $\CSpec$ subsumes any pairwise
compatibility check that $\mathcal{V}$ performs.
\end{proposition}

\chapter{GOFAI Natural and Creative Language}
\label{ch:gofai-nl}

A central claim of Judgment Geometry is that structured mathematical
theory can \emph{replace} learned statistical models for certain
language tasks.  This chapter develops the claim in full: we show how
the Compositional Theory of Theories (CTT), instantiated with $\sim\!450$
theory slots across ten linguistic strata, drives a chatbot and a
poetry engine that use \emph{no neural model whatsoever}---zero
parameters, zero training data, zero GPU.

%% ─── 28.1 ──────────────────────────────────────────────────────
\section{The GOFAI thesis}
\label{sec:gofai-thesis}

\begin{quote}
\emph{If you have enough theory, you don't need an LLM.}
\end{quote}

The CTT of Part~III provides:
\begin{enumerate}[label=(\arabic*)]
  \item \textbf{450 theory slots} spanning phonology through aesthetics,
        each contributing constraints and generation rules.
  \item \textbf{The $\Grade$ semiring} for combining theory verdicts:
        $\otG$ conjoins independent requirements; $\opG$ selects
        among competing alternatives.
  \item \textbf{The $\FMLM$ decoder}, which fills masked positions
        by iterative fixpoint computation over the full theory
        ensemble.
  \item \textbf{The descent algorithm}, which repairs
        disharmonious output until the global grade crosses a
        threshold.
\end{enumerate}

The result is a system whose every output is \emph{explainable}
(each quality grade traces to a named theory), \emph{controllable}
(change one theory slot, change one aspect of output), and
\emph{compositional} (theories compose via the semiring without
interference).

%% ─── 28.2 ──────────────────────────────────────────────────────
\section{The GOFAI chatbot architecture}
\label{sec:gofai-chatbot}

The chatbot implements a classic pipeline, but every stage is
graded by a JuGeo theory slot.

\subsection{Natural language understanding}
\label{ssec:nlu}

Input processing follows a five-stage NLU pipeline:
\begin{enumerate}
  \item \textbf{Tokenization and POS tagging}\; ($\morph$ stratum).
        Rule-based tokenizer with lexicon lookup; POS tags assigned by
        suffix rules and disambiguation via the highest-grade parse.
  \item \textbf{Syntactic parsing}\; ($\synt$ stratum).
        A slot-based parser identifies clause boundaries, verb frames,
        and modifier attachment.  Each parse receives a grade
        $g_{\synt}$ from theories such as \textsf{ClauseCompleteness}
        and \textsf{InversionLicensing}.
  \item \textbf{Semantic analysis}\; ($\sem$ stratum).
        Frame-semantic analysis using FrameNet frames: the input is
        mapped to a frame with filled roles.  Theories
        \textsf{FrameCoherence} and \textsf{EntityBinding} grade the
        result.
  \item \textbf{Pragmatic interpretation}\; ($\prag$ stratum).
        Speech-act classification (request, assertion, question) and
        Gricean maxim checking.  The theory \textsf{SpeechActFelicity}
        grades the pragmatic parse.
  \item \textbf{Discourse integration}\; ($\disc$ stratum).
        Topic tracking and anaphora resolution update the discourse
        state.  Theories \textsf{TopicCoherence} and
        \textsf{AnaphoraResolution} grade the integration.
\end{enumerate}

The overall NLU grade is the $\otG$-product across strata:
\begin{equation}
\label{eq:nlu-grade}
  g_{\mathrm{NLU}} = g_{\morph} \otG g_{\synt} \otG g_{\sem}
    \otG g_{\prag} \otG g_{\disc}.
\end{equation}

\subsection{Intent classification}
\label{ssec:intent}

The dialogue manager classifies user intent via a pattern-based
system graded by the $\Grade$ semiring.  Each candidate intent
$I \in \{\textsc{PoemRequest}, \textsc{StyleRequest},
\textsc{Feedback}, \textsc{Revision}, \textsc{Question},
\textsc{Greeting}, \textsc{Help}, \textsc{Thanks}\}$ has a set of
indicator patterns; firing patterns contribute grade via~$\opG$
(log-sum-exp), and the highest-grade intent wins:
\begin{equation}
\label{eq:intent-grade}
  I^* = \arg\max_I \; g(I \mid \text{input}).
\end{equation}
No neural classifier is involved: the grade semiring \emph{is} the
classifier.

\subsection{Response generation}
\label{ssec:response-gen}

Given the classified intent $I^*$ and discourse state, the response
pipeline proceeds:
\begin{enumerate}
  \item \textbf{Topic model $\to$ frame selection.}\;
        The current topic (maintained by the $\disc$ stratum) selects
        a FrameNet frame; frame elements become template slots.
  \item \textbf{Template instantiation.}\;
        A response template is chosen from a bank indexed by intent
        and frame; slots are filled from the discourse context and
        knowledge base.
  \item \textbf{Surface realization.}\;
        The filled template is linearized, applying morphological
        rules ($\morph$ stratum) and phonological smoothing ($\phon$
        stratum).
  \item \textbf{$\FMLM$ refinement.}\;
        Any remaining masked positions are filled by the fixpoint
        decoder, which scores candidates against the full theory
        ensemble.
\end{enumerate}

\subsection{Conversation management}
\label{ssec:conversation}

Long-range coherence is maintained by the \emph{discourse monad}
$\mathrm{DMon}$, which tracks:
\begin{itemize}
  \item \textbf{Topic stack}\; ($\disc$ stratum): active topics,
        topic shifts, and topic resumption.
  \item \textbf{Reference resolution}\; ($\sem$ stratum): entity
        mentions and their antecedents.
  \item \textbf{Turn history}\; ($\prag$ stratum): user preferences,
        feedback history, and affective state.
  \item \textbf{User model}: accumulated preferences for poetic form,
        mood, and style, updated after each interaction.
\end{itemize}

The monadic structure ensures that state updates compose cleanly
across turns: each turn is a Kleisli arrow
$\mathrm{Turn} \to \mathrm{DMon}(\mathrm{Response})$.

%% ─── 28.3 ──────────────────────────────────────────────────────
\section{The GOFAI poetry engine}
\label{sec:gofai-poetry}

The poetry engine is the showcase application of the GOFAI approach.
It generates metrically, rhymically, and semantically controlled
verse using no neural model.

\subsection{Form selection}
\label{ssec:form-select}

Each poetic form is a type $F$ in $\PStr$ with specific structural
constraints:
\begin{center}
\begin{tabular}{lll}
\textbf{Form} & \textbf{Constraint type} & \textbf{Theory} \\
\hline
Sonnet & 14 lines, volta at 8 or 12, iambic pentameter
       & \textsf{FormCompliance} \\
Villanelle & 19 lines, two refrains, $\mathrm{ABA}$ tercets
           & \textsf{FormCompliance} \\
Haiku & 5-7-5 syllables, seasonal word (\emph{kigo})
      & \textsf{SyllableCount} \\
Ghazal & couplets with \emph{radif} and \emph{qafia}
       & \textsf{RhymeScheme} \\
Terza rima & $\mathrm{ABA\;BCB\;CDC\;\dots}$ interlocking tercets
           & \textsf{RhymeScheme} \\
Free verse & no fixed meter; scored on rhythm variety
           & \textsf{MeterCompliance}
\end{tabular}
\end{center}
The user may also specify a \emph{poet style} (Keats, Dickinson,
Whitman, Shakespeare, Neruda), which modulates diction, syntax, and
thematic profiles via a style presheaf over the stratal tower.

\subsection{Content generation}
\label{ssec:content-gen}

Content is generated by pure knowledge-base lookup:
\begin{enumerate}
  \item \textbf{Topic $\to$ FrameNet frames.}\;
        The topic word activates one or more FrameNet frames
        (e.g., \emph{love} activates \textsc{Personal\_relationship},
        \textsc{Emotion\_directed}, \textsc{Desiring}).
  \item \textbf{Frame $\to$ lexical units.}\;
        Each frame provides lexical units; these are expanded via
        WordNet synonymy, hypernymy, and semantic similarity.
  \item \textbf{Candidate word pool.}\;
        The union of lexical units and WordNet expansions forms
        the candidate pool, from which the $\FMLM$ decoder draws.
  \item \textbf{Proof strategies.}\;
        Eight proof strategies---\textsf{BaseFact}, \textsf{Causal},
        \textsf{Metaphor}, \textsf{Temporal}, \textsf{Negation},
        \textsf{Relational}, \textsf{Hypernym},
        \textsf{Presupposition}---each produce an $\HPLF$ term and
        a template with bindings.  These provide the propositional
        skeleton of each line.
\end{enumerate}

No neural model is consulted at any point: content derives from
FrameNet, WordNet, CMUDict, and the theory ensemble.

\subsection{Quality control}
\label{ssec:quality}

Each candidate poem is scored by the $\Grade$ semiring along four
primary axes, each backed by specific theory slots:

\medskip\noindent\textbf{Meter.}\;
CMUDict provides stress patterns (sequences of 0, 1, 2 for
unstressed, primary, secondary stress).  The theory
\textsf{MeterCompliance} computes the edit distance between the
observed stress pattern and the target pattern (e.g., iambic
pentameter = $01\,01\,01\,01\,01$), normalized to $[-\infty, 0]$:
\[
  g_{\mathrm{meter}} = -\frac{d_{\mathrm{edit}}(\sigma_{\mathrm{obs}},
    \sigma_{\mathrm{target}})}{\lvert \sigma_{\mathrm{target}} \rvert}.
\]

\medskip\noindent\textbf{Rhyme.}\;
CMUDict phoneme sequences enable perfect rhyme (identical final
stressed vowel and coda), near rhyme (identical vowel, different
coda), and slant rhyme (similar vowel).  The theory
\textsf{RhymeScheme} checks each line-ending pair against the
target scheme:
\[
  g_{\mathrm{rhyme}} = \frac{1}{|\text{pairs}|}\sum_{\text{pairs}}
    g_{\mathrm{match}}(w_i, w_j).
\]

\medskip\noindent\textbf{Imagery.}\;
WordNet Wu--Palmer similarity between successive content words
measures semantic diversity; excessively high similarity (clich\'e)
or excessively low similarity (incoherence) is penalized:
\[
  g_{\mathrm{imagery}} = -\bigl\lvert
    \mathrm{sim}(w_i, w_j) - \tau_{\mathrm{target}} \bigr\rvert,
\]
where $\tau_{\mathrm{target}} \approx 0.4$ balances novelty and
coherence.

\medskip\noindent\textbf{Emotion.}\;
Affect lexicons (valence, arousal, dominance) map each word to an
emotional vector.  The theory \textsf{RasaConsistency} checks that
the poem's emotional trajectory matches the target
mood (\emph{rasa}):
\[
  g_{\mathrm{emotion}} = -\left\lVert
    \vec{e}_{\mathrm{poem}} - \vec{e}_{\mathrm{target}}
  \right\rVert.
\]

The total poem grade is:
\begin{equation}
\label{eq:poem-grade}
  g_{\mathrm{poem}} = g_{\mathrm{meter}} \otG g_{\mathrm{rhyme}}
    \otG g_{\mathrm{imagery}} \otG g_{\mathrm{emotion}}
    \otG g_{\mathrm{style}} \otG g_{\mathrm{form}}.
\end{equation}

\subsection{Beam search as proof search}
\label{ssec:beam-search}

The generation engine maintains a beam of $k$ candidate poems,
scored by their total grade~\eqref{eq:poem-grade}.  At each
step (line or stanza), the $\FMLM$ decoder produces $m$ candidate
continuations for each beam element; the top-$k$ by grade survive.
This is formally identical to \emph{proof search} in the type
theory of Part~III: each candidate poem is a partial proof term,
each continuation extends the proof, and the grade is the
proof's quality.

\begin{proposition}
\label{prop:beam-convergence}
If the per-step grade is bounded below by $g_{\min} > \gzero$ and
the form imposes a finite length~$L$, then beam search terminates
in at most $L$ steps with at least one candidate of grade
$\geq g_{\min}^L$ (where exponentiation is in the $\otG$-monoid,
i.e., $L$-fold $\otG$-product).
\end{proposition}

%% ─── 28.4 ──────────────────────────────────────────────────────
\section{Advantages of GOFAI + JuGeo}
\label{sec:gofai-advantages}

We summarize the structural advantages of the theory-driven approach
over LLM-based generation.

\begin{description}
  \item[Interpretability.]
    Every component of $g_{\mathrm{poem}}$ traces to a named theory
    slot.  A user can ask ``why is the meter grade low?'' and receive
    a specific answer (``line~7 has a trochaic inversion at
    position~3'').  An LLM cannot provide this.
  \item[Controllability.]
    Changing one theory slot changes exactly one aspect of the
    output.  To make the poem more emotionally intense, raise the
    weight of \textsf{RasaConsistency}; to allow more metrical
    freedom, lower the weight of \textsf{MeterCompliance}.
    In an LLM, such fine-grained control is unavailable.
  \item[Compositionality.]
    Theories compose via the $\Grade$ semiring: adding a new
    theory (e.g., \textsf{Alliteration}) does not interfere with
    existing ones.  This is the \emph{open world property} of CTT.
  \item[No training data.]
    The system runs on pure theory: lexicons (WordNet, CMUDict,
    FrameNet), affect dictionaries, and hand-authored constraints.
    No corpus of poems is needed.
  \item[No hallucination.]
    Every factual claim in the poem traces to a knowledge-base
    entry (WordNet gloss, FrameNet frame, or Wikipedia fact).
    The system cannot invent false facts because it has no
    generative model to hallucinate from.
\end{description}

%% ─── 28.5 ──────────────────────────────────────────────────────
\section{Disadvantages and mitigations}
\label{sec:gofai-disadvantages}

\begin{description}
  \item[Fluency.]
    GOFAI-generated text can sound stilted.  Mitigation: large word
    banks ($> 150{,}000$ words from CMUDict), frame-based templates
    from FrameNet, and the $\FMLM$ decoder's multi-axis scoring
    (phonological, semantic, register, diversity) smooth the surface.
  \item[Novelty.]
    Knowledge-base lookup risks producing predictable combinations.
    Mitigation: the \emph{conceptual blending} type constructor
    $\mathrm{Blend}(F_1, F_2)$ of Part~III merges two FrameNet frames
    to produce novel juxtapositions (e.g., blending
    \textsc{Commerce\_sell} with \textsc{Death} for a war poem).
  \item[Scale.]
    Full beam search over 450 theories is expensive.
    Mitigation: theories are lazily evaluated (only active
    theories fire), and the $\FMLM$ decoder caches lexical lookups
    across lines.  In practice, generation of a 14-line sonnet
    takes $<$\,2 seconds on commodity hardware.
\end{description}

%% ─── 28.6 ──────────────────────────────────────────────────────
\section{The hybrid approach}
\label{sec:hybrid}

The GOFAI and LLM approaches are not mutually exclusive; JuGeo
provides the framework for combining them.

\begin{construction}[Hybrid Generation Pipeline]
\label{con:hybrid}
\begin{enumerate}
  \item \textbf{GOFAI generates the structural skeleton.}\;
        The theory ensemble produces a metrically and rhymically
        correct draft with filled semantic roles, graded by
        $\GEval$.  This draft is a \emph{section} in $\PStr$ with
        $\check{H}^1 = 0$ (it glues) but potentially low surface
        fluency.
  \item \textbf{LLM polishes the surface.}\;
        The draft is passed to an LLM with the instruction ``improve
        fluency while preserving meter, rhyme, and meaning.''
        The LLM's output is a new candidate section.
  \item \textbf{JuGeo verifies the result.}\;
        The full theory ensemble re-grades the polished version.
        If $\GEval$ improves and $\check{H}^1$ remains zero, the
        polished version is accepted; otherwise the GOFAI draft is
        preferred.
\end{enumerate}
\end{construction}

This hybrid pipeline gives the best of both worlds: the structural
guarantees of GOFAI and the surface fluency of an LLM, with JuGeo
serving as the arbiter.  The graded evaluation ensures that the
LLM cannot silently break structural invariants---any fluency
improvement that damages meter, rhyme, or semantic coherence will
be caught by the theory ensemble and rejected.

\chapter{Software Generation at Scale}
\label{ch:software-scale}

We now turn from individual-project verification to
\emph{industrial-scale} applications of Judgment Geometry:
continuous integration, multi-repository coherence, specification
mining, refactoring, security analysis, and performance optimization.

%% ─── 29.1 ──────────────────────────────────────────────────────
\section{Continuous verification}
\label{sec:continuous-verif}

Modern software development is organized around
continuous-integration (CI) pipelines.  JuGeo integrates naturally
into this workflow.

\begin{construction}[JuGeo CI Layer]
\label{con:ci-layer}
Augment a standard CI/CD pipeline with a JuGeo verification stage:
\begin{enumerate}
  \item \textbf{Every commit $c$ induces a judgment.}\;
        The diff of $c$ modifies sections on one or more open sets
        $U_i$ in the project site $\STop$.  JuGeo recomputes the
        affected entries of the \v{C}ech complex.
  \item \textbf{Every pull request triggers descent.}\;
        The PR's changeset defines a candidate section~$s'$;
        JuGeo computes $\delta^0 s'$ and reports any new
        obstructions introduced by the change.
  \item \textbf{$\CSpec$ catches integration bugs.}\;
        Standard CI runs unit tests ($H^0$ check).  JuGeo adds the
        $H^1$ check: do the changed modules still glue with their
        neighbors?
  \item \textbf{$\GEval$ tracks quality over time.}\;
        The grade $g(s)$ is computed for every merge to the main
        branch, producing a \emph{quality time series}
        $\{g(s^{(t)})\}_{t \geq 0}$.  Regressions are flagged when
        $g(s^{(t+1)}) < g(s^{(t)}) - \epsilon$ for a configurable
        tolerance~$\epsilon$.
\end{enumerate}
\end{construction}

\begin{remark}
The JuGeo CI layer is \emph{incremental}: only the affected
part of the \v{C}ech complex is recomputed.  If a commit touches
modules $M_i$ and~$M_j$, only the overlap cochains involving $M_i$
and~$M_j$ are re-evaluated.  This makes the computation
proportional to the size of the diff, not the size of the project.
\end{remark}

%% ─── 29.2 ──────────────────────────────────────────────────────
\section{Multi-repository coherence}
\label{sec:multi-repo}

Large organizations operate across monorepos or constellations of
microservice repositories.  JuGeo scales naturally to this setting.

\begin{definition}[Multi-Repository Site]
\label{def:multi-repo-site}
Let $R_1, \dots, R_n$ be repositories.  Define the
\emph{multi-repository site} $\STop_{\mathrm{multi}}$:
\begin{itemize}
  \item Objects: the modules of each $R_i$, plus inter-repository
        API surfaces.
  \item Morphisms: intra-repo dependency + inter-repo API calls.
  \item Topology: a family covers an API surface~$A$ if it
        includes both the provider and all consumers of~$A$.
\end{itemize}
\end{definition}

Under this site:
\begin{itemize}
  \item \textbf{API contracts are $\Cech^1$ conditions.}\;
        The overlap $U_i \cap U_j$ between a service and its
        client is the API contract (request/response schema,
        error codes, latency bounds).  A class
        $[\alpha] \in \check{H}^1$ with $[\alpha]\neq 0$ signals a
        \emph{contract violation}.
  \item \textbf{Breaking changes are $H^1 \neq 0$.}\;
        A breaking API change is precisely a modification that
        introduces a non-trivial class in $\check{H}^1$: the
        provider and consumer disagree on the overlap.
  \item \textbf{Versioning is cover refinement.}\;
        Introducing a new API version $v2$ alongside $v1$ refines
        the cover: consumers of $v1$ and $v2$ live in different
        open sets, and the overlap conditions split accordingly.
        This is sheaf-theoretically natural---it corresponds to
        passing to a finer topology.
\end{itemize}

\begin{proposition}
\label{prop:semver}
Under the multi-repository site, semantic versioning
(major/minor/patch) admits the following cohomological
interpretation:
\begin{enumerate}[label=(\roman*)]
  \item A \emph{patch} change preserves all sections
        ($\check{H}^1$ unchanged).
  \item A \emph{minor} change adds new sections but
        preserves existing ones
        ($\check{H}^1$ can only shrink).
  \item A \emph{major} change may introduce new classes in
        $\check{H}^1$ (existing sections may fail to glue).
\end{enumerate}
\end{proposition}

%% ─── 29.3 ──────────────────────────────────────────────────────
\section{Specification mining}
\label{sec:spec-mining}

A key practical obstacle to verification is that most code lacks
formal specifications.  JuGeo can \emph{infer} specifications from
existing code and tests.

\begin{construction}[Specification Mining via $\PStr$]
\label{con:spec-mining}
Given a module~$U$ with code~$E_U$ and test suite~$T_U$:
\begin{enumerate}
  \item \textbf{Mine candidate types.}\;
        Analyze $E_U$ to extract function signatures, invariants,
        and pre/post-conditions.  Each candidate is a potential
        type $A_U$ in~$\PStr$.
  \item \textbf{Tests as evidence.}\;
        Each passing test $t \in T_U$ is an evidence term
        $e_t : A_U$ in the mined type.  The grade
        $g(e_t : A_U)$ measures how well the test exercises the
        specification.
  \item \textbf{Cross-module consistency.}\;
        Compute $\check{H}^1$ for the mined specifications:
        if module $M_i$'s mined spec says ``\texttt{get\_user}
        returns a string'' but $M_j$'s tests treat the return
        value as a dictionary, $\check{H}^1 \neq 0$ flags the
        inconsistency.
  \item \textbf{Refinement.}\;
        Use the obstructions to refine the mined specifications
        until $\check{H}^1 = 0$.  The result is a consistent,
        machine-inferred specification for the entire project.
\end{enumerate}
\end{construction}

\begin{remark}
Specification mining is dual to code generation: generation starts
from $A$ and produces~$E$; mining starts from $E$ and infers~$A$.
In both cases, $\CSpec$ enforces global consistency.
\end{remark}

%% ─── 29.4 ──────────────────────────────────────────────────────
\section{Refactoring as descent}
\label{sec:refactoring}

Large-scale refactoring---renaming, module extraction, API
migration---is one of the most error-prone activities in software
engineering.  Under JuGeo, refactoring is \emph{exactly} descent.

\begin{proposition}[Refactoring = Descent]
\label{prop:refactoring}
Let $s^{(0)}$ be the current codebase (a section of $\PStr$) and
let $A'$ be the target specification (post-refactoring types and
interfaces).  Refactoring from $s^{(0)}$ to a section $s^{(*)}$
satisfying $A'$ is the descent problem for the presheaf
$\PStr_{A'}$ with initial candidate~$s^{(0)}$.  Specifically:
\begin{enumerate}
  \item The current code $s^{(0)}$ does not satisfy $A'$ on overlaps:
        $\check{H}^1(\STop, \PStr_{A'}) \neq 0$ with $s^{(0)}$ as
        candidate.
  \item Each refactoring step modifies $s^{(t)}$ on some $U_i$,
        producing $s^{(t+1)}$.
  \item Descent converges when $\check{H}^1 = 0$ relative to~$A'$.
\end{enumerate}
\end{proposition}

The $\GEval$ component tracks refactoring progress:
$g(s^{(t)})$ increases monotonically if each step resolves at least
one obstruction.  This provides a \emph{progress metric} for
refactoring projects that is more informative than ``percentage of
files changed.''

%% ─── 29.5 ──────────────────────────────────────────────────────
\section{Security analysis via cohomology}
\label{sec:security}

Security vulnerabilities often arise from \emph{non-local}
interactions: data flows from a trusted source through an untrusted
transformation to a sensitive sink.  This is a cohomological
phenomenon.

\begin{construction}[Security Presheaf]
\label{con:security-presheaf}
Define a presheaf $\mathcal{S}$ on the call-graph site of a program:
\begin{itemize}
  \item To each function $f$, assign the set of
        \emph{security labels} (e.g., public, confidential, secret)
        of its inputs and outputs.
  \item To each call edge $f \to g$, assign the restriction map
        that propagates labels through the call.
\end{itemize}
The sheaf condition requires that labels are consistent along every
path from source to sink.  Failure of the sheaf condition---a
non-trivial class in $\check{H}^1(\mathcal{G}, \mathcal{S})$---is
an \emph{information leak}: data classified ``secret'' reaches a
function labeled ``public'' via an intermediate path that fails to
enforce the label.
\end{construction}

\begin{example}
A web application has a function \texttt{get\_password} labeled
\texttt{secret} and a logging function \texttt{log} labeled
\texttt{public}.  An intermediate function \texttt{format\_user}
calls both.  The overlap cochain records that \texttt{format\_user}
must simultaneously handle \texttt{secret} input and produce
\texttt{public} output---a non-trivial $H^1$ class, flagging a
potential leak path.
\end{example}

Patching the vulnerability is descent: modify the code (repair the
section) until the security labels glue consistently and
$\check{H}^1 = 0$.

%% ─── 29.6 ──────────────────────────────────────────────────────
\section{Performance optimization}
\label{sec:perf-opt}

Performance can be treated as an additional stratum in the JuGeo
framework.

\begin{construction}[Performance Stratum]
\label{con:perf-stratum}
Augment the project site with a \emph{performance stratum}
$\sigma_{\mathrm{perf}}$:
\begin{itemize}
  \item To each function~$f$, assign a performance profile
        (time complexity, memory usage, cache behavior).
  \item The grade $g_{\mathrm{perf}}(f)$ measures how well $f$'s
        observed performance matches its target.
  \item \emph{Hotspots} are obstructions: functions where
        $g_{\mathrm{perf}} \ll \gone$.
  \item \emph{Optimization} is descent: iteratively improve the
        lowest-grade functions until
        $g_{\mathrm{perf}} \geq g_{\min}$.
\end{itemize}
\end{construction}

Cross-stratum descent is particularly powerful here: a performance
optimization that changes a function's interface (e.g., switching
from synchronous to asynchronous I/O) creates an obstruction in
the functional stratum.  JuGeo's cross-stratum descent detects this
and ensures that all callers are updated accordingly.

\begin{remark}
The performance stratum illustrates the \emph{extensibility} of
JuGeo: adding a new stratum requires only defining the appropriate
theory slots and grade functions.  The descent algorithm, the
cohomological machinery, and the $\Grade$ semiring are reused
unchanged.  This is the open-world property of CTT applied to
software engineering.
\end{remark}

\chapter{Future Directions}
\label{ch:future}

We close the monograph with a survey of open problems, speculative
applications, and directions for future research.  Some of these are
natural extensions of the existing framework; others require
fundamentally new mathematics.

%% ─── 30.1 ──────────────────────────────────────────────────────
\section{Higher-categorical Judgment Geometry}
\label{sec:higher-cat}

Throughout this monograph, $\STop$ has been a 1-topos: a category
of sheaves on a Grothendieck site.  The natural generalization
replaces 1-categories with $(\infty,1)$-categories.

\begin{itemize}
  \item \textbf{$\STop$ becomes an $\infty$-topos} in the sense of
        Lurie~\cite{lurie2009}.  The judgment presheaf $\PStr$
        becomes an $\infty$-presheaf valued in $\infty$-groupoids
        rather than sets; this captures not just whether two
        local sections agree on an overlap, but all \emph{higher
        coherences} (homotopies between homotopies, \emph{ad
        infinitum}).
  \item \textbf{$\PStr$ becomes homotopy type theory.}\;
        Types are no longer mere propositions or sets but
        \emph{spaces}; the identity type $\GId_g(a,b)$ of Part~III
        generalizes to a graded path space, and higher identity
        types capture graded homotopies.
  \item \textbf{$\CSpec$ becomes derived $\infty$-cohomology.}\;
        The \v{C}ech complex is replaced by the full derived
        functor of global sections in the $\infty$-topos.  The
        obstruction theory gains access to the full spectrum of
        cohomological invariants, including $H^n$ for all $n \geq 0$.
  \item \textbf{$\GEval$ becomes an $\infty$-graded structure.}\;
        The $\Grade$ semiring is replaced by a
        \emph{graded $\infty$-ring spectrum}, in which grades
        themselves form a homotopy-coherent algebraic structure.
\end{itemize}

\begin{remark}
The motivation for higher-categorical JuGeo is not merely
generality.  In software verification, higher coherences arise
naturally: a refactoring $r_1$ and a refactoring $r_2$ may each
preserve correctness, but their composition $r_2 \circ r_1$ may
not---and tracking \emph{why} requires a homotopy between the
two paths.  In linguistics, higher coherences appear in poetic
ambiguity: a word may have two readings, each reading two
sub-readings, and so forth.  The $\infty$-categorical framework
captures this tower of ambiguity natively.
\end{remark}

The key open problem is:

\begin{quote}
\emph{Does the descent algorithm of Part~II generalize to the
$\infty$-categorical setting with convergence guarantees?}
\end{quote}

We conjecture that it does, provided the $\infty$-graded structure
satisfies an appropriate finiteness condition (e.g., that the
homotopy groups of the grade spectrum stabilize).

%% ─── 30.2 ──────────────────────────────────────────────────────
\section{Quantum Judgment Geometry}
\label{sec:quantum}

Quantum computing introduces programs whose correctness cannot be
verified by classical testing (measuring a quantum state collapses
it).  JuGeo offers an alternative.

\begin{itemize}
  \item \textbf{The site} is the category of quantum circuits,
        with morphisms given by circuit composition and ancilla
        management.
  \item \textbf{Types} are quantum specifications: unitary
        conditions ($U^\dagger U = I$), entanglement invariants,
        and decoherence bounds.
  \item \textbf{Cohomology} captures topological invariants of
        quantum states---analogous to topological quantum
        computation, where the robustness of a quantum gate is
        measured by a topological invariant.
  \item \textbf{Grades} are fidelity measures: the grade
        $g(\rho, \sigma) = \log F(\rho, \sigma)$ where $F$ is the
        quantum fidelity between the ideal state $\sigma$ and the
        actual state~$\rho$.
\end{itemize}

\begin{remark}
The connection between JuGeo and topological quantum computing
is tantalizing: both frameworks use cohomological invariants to
measure the ``robustness'' of a structure.  Whether this analogy
can be made precise---whether quantum error correction codes
arise as descent data in a JuGeo site over quantum circuits---is
an open question.
\end{remark}

%% ─── 30.3 ──────────────────────────────────────────────────────
\section{Biological Judgment Geometry}
\label{sec:biological}

Biological macromolecules can be viewed as ``programs'' whose
``execution'' is physical.  JuGeo provides a verification framework.

\begin{itemize}
  \item \textbf{Codons as morphemes.}\;
        The genetic code maps nucleotide triplets to amino acids,
        exactly as morphological rules map morphemes to meanings.
        The $\morph$ stratum of CTT applies directly.
  \item \textbf{Protein folding as descent.}\;
        A protein's amino acid sequence is a ``local'' description;
        the folded 3D structure is the ``global section'' that the
        sequence must descend to.  Misfolding is $H^1 \neq 0$: the
        local energetic preferences of individual residues fail to
        glue into a globally consistent structure.
  \item \textbf{Mutations as obstructions.}\;
        A point mutation modifies the section on one open set
        (one residue's local environment).  Whether the mutant
        protein still folds correctly is whether the modified
        section still glues---a cohomological question.
  \item \textbf{Fitness as grade.}\;
        Darwinian fitness provides a natural $\GEval$: the grade
        of an organism is its reproductive success, and evolution
        is descent toward higher grades.
\end{itemize}

%% ─── 30.4 ──────────────────────────────────────────────────────
\section{Musical Judgment Geometry: DAW Projects as Programs}
\label{sec:musical}

A digital audio workstation (DAW) project file---a Renoise
\texttt{.xrns}, an Ableton Live \texttt{.als}, an FL~Studio
\texttt{.flp}---is a \emph{complete program}.  It contains tracks
(modules), patterns (functions), notes and effect commands
(instructions), signal-flow routing (data flow), automation curves
(continuous parameter functions), a mixer dependency graph (sends,
returns, sidechains), references to external plugins (libraries),
and a master bus (entry point).  The analogy is not metaphorical:
these project files admit exactly the same 4-tuple structure
$\JJ = (\STop, \PStr, \CSpec, \GEval)$ that governs software
verification.  In this section we develop that structure in detail,
using tracker-style DAWs (especially Renoise) as the running
example while noting how each concept maps to timeline-based
environments such as Ableton Live and FL~Studio.

% ────────────────────────────────────────────────────────────────
\subsection{The site \texorpdfstring{$\STop$}{S} for a DAW project}
\label{ssec:musical-site}

The site is the temporal-structural hierarchy of the project.

\begin{definition}[Musical site]
\label{def:musical-site}
Let $\mathcal{C}_{\mathrm{DAW}}$ be the category whose objects are
the structural units of the project at every scale:
\begin{enumerate}
  \item \textbf{Sample level.}\;
        Individual audio samples at the sample rate
        (e.g.\ 44\,100\,Hz or 48\,000\,Hz).
  \item \textbf{Tick level.}\;
        Tracker ticks or MIDI ticks---the finest rhythmic
        resolution.  In Renoise the default is 12~ticks per line;
        at 8~lines per beat this gives 96~ticks per beat.
  \item \textbf{Line / beat level.}\;
        Individual rows in a Renoise pattern (``lines''),
        corresponding to sixteenth-note or other subdivisions,
        or beats in a timeline DAW.
  \item \textbf{Pattern level.}\;
        Renoise's fundamental compositional unit: a rectangular
        grid of note columns and effect columns, of configurable
        length.  In Ableton this corresponds to a clip; in
        FL~Studio, a pattern clip.
  \item \textbf{Track level.}\;
        Instrument tracks, group tracks, send tracks, and the
        master track.
  \item \textbf{Song level.}\;
        The pattern sequence (Renoise's ``order list''), or the
        arrangement view timeline.
  \item \textbf{Project level.}\;
        The complete \texttt{.xrns}/\texttt{.als}/\texttt{.flp}
        file, including all referenced samples, instruments,
        plugin states, and configuration.
\end{enumerate}
Morphisms are \emph{restriction} maps: a pattern restricts to a
line, a line restricts to a tick, a track restricts to its content
within a single pattern, etc.  Composition of restrictions is
associative; identity restrictions exist at every level.
\end{definition}

The Grothendieck topology on $\mathcal{C}_{\mathrm{DAW}}$ is
defined by covering sieves:

\begin{itemize}
  \item A \emph{cover of the song} is its pattern sequence---the
        ordered collection of pattern instances that tile the
        timeline.
  \item A \emph{cover of a pattern} is its track decomposition---
        the collection of per-track columns within that pattern.
  \item The \emph{overlap} between adjacent patterns in the
        sequence captures crossfade regions, automation
        continuity at pattern boundaries, and sustained notes
        that cross a pattern break.
\end{itemize}

In Renoise, the \textbf{pattern matrix}---the grid of pattern
slots vs.\ tracks---\emph{is} the covering sieve, made visible as
a GUI element.  Each cell $(p, t)$ in the matrix is an open set
$U_{p,t}$, and the pattern matrix is the nerve of the cover.

\begin{definition}[Musical presheaf]
\label{def:musical-presheaf}
The \emph{structure presheaf} $\mathcal{F}$ on
$\mathcal{C}_{\mathrm{DAW}}$ assigns to each open set $U$ the
musical data active on $U$:
\begin{itemize}
  \item At a line: note values, velocities, instrument indices,
        panning, volume, and effect-column commands
        (e.g.\ \texttt{0Exx} retrigger, \texttt{0Bxx} pattern
        break, \texttt{0900} sample offset).
  \item At a track: the DSP effect chain, send amounts, group
        routing.
  \item At an instrument: sample mappings, modulation sets
        (envelopes, LFOs), macro assignments, plugin state.
  \item At a pattern: automation envelopes for any device
        parameter.
\end{itemize}
Restriction maps are the natural projections:
$\mathcal{F}(U) \to \mathcal{F}(V)$ for $V \hookrightarrow U$
restricts the data to the sub-region.
\end{definition}

\begin{remark}[Renoise effect commands as morphisms]
Pattern effect commands such as \texttt{0Bxx} (pattern break) and
\texttt{0Fxx} (set speed/LPB) are morphisms \emph{between} pattern
objects: they alter the flow of control in the pattern sequence,
just as jump instructions alter control flow in a program.
Instrument macros---Renoise's 8 assignable knobs that each control
an arbitrary set of device parameters---are natural transformations
between sub-presheaves, coordinating multiple restriction maps
through a single control surface.
\end{remark}

% ────────────────────────────────────────────────────────────────
\subsection{The proof structure \texorpdfstring{$\PStr$}{P} for
  music}
\label{ssec:musical-proof}

The \emph{type} of a DAW project is its musical specification---the
constraints that the finished piece must satisfy.

\begin{definition}[Musical type]
\label{def:musical-type}
A \emph{musical type} $\tau$ is a tuple of constraints drawn from:
\begin{itemize}
  \item \textbf{Global parameters.}\;
        BPM range, time signature, key, scale/mode.
  \item \textbf{Arrangement structure.}\;
        An ordered list of section labels---e.g.\
        $[\mathsf{intro}, \mathsf{verse}, \mathsf{chorus},
         \mathsf{bridge}, \mathsf{chorus}, \mathsf{outro}]$---with
        target durations.
  \item \textbf{Harmonic rules.}\;
        Chord progression grammars (e.g.\ $\mathrm{ii}$--$\mathrm{V}$--$\mathrm{I}$,
        circle-of-fifths motion), modulation constraints,
        voice-leading rules.
  \item \textbf{Rhythmic constraints.}\;
        Groove template, swing percentage, quantization grid.
  \item \textbf{Timbral constraints.}\;
        Frequency band allocation per track group, stereo-field
        assignments (e.g.\ bass and kick centered, hats and
        percussion panned), dynamic range targets.
  \item \textbf{Reference targets.}\;
        A reference track with measured LUFS loudness,
        spectral profile, and stereo width---the ``type'' the
        final master should inhabit.
\end{itemize}
\end{definition}

The \emph{evidence} inhabiting this type is the actual musical
content: the patterns, notes, automation curves, DSP chains, and
mixer settings that constitute the project.  The \emph{typing
derivation} is the compositional structure that connects each piece
of evidence to its constraint:

\begin{itemize}
  \item Each pattern's content derives from the arrangement type
        for that section.
  \item Each note in a pattern derives from the harmonic type
        active at that moment.
  \item Each DSP device in a chain derives from the timbral type
        for that track's frequency/stereo allocation.
  \item Each automation curve derives from the dynamic contour
        specified by the arrangement arc.
\end{itemize}

\begin{definition}[Renoise-specific types]
\label{def:renoise-types}
We define the following type constructors for Renoise projects:
\begin{align*}
  &\mathsf{PatternType}(\ell, n_t, \mathit{lpb})
    &&\text{--- pattern of length $\ell$ lines, $n_t$ tracks,
           $\mathit{lpb}$ lines per beat,} \\
  &\mathsf{InstrumentType}(S, M, K)
    &&\text{--- instrument with sample set $S$, modulation sets $M$,
           macro map $K$,} \\
  &\mathsf{EffectChainType}(D_1, \ldots, D_k)
    &&\text{--- chain of DSP devices $D_i$ in signal-flow order,} \\
  &\mathsf{AutomationType}(p, e)
    &&\text{--- automation envelope $e$ for device parameter $p$.}
\end{align*}
\end{definition}

A well-typed project is one where every piece of musical content
has a derivation from the specification.  An untyped passage---a
note outside the scale, a pattern with no arrangement role, an
unrouted send---is a type error.

% ────────────────────────────────────────────────────────────────
\subsection{The cohomological spectrum \texorpdfstring{$\CSpec$}{H}
  for music}
\label{ssec:musical-cohomology}

Cohomology detects precisely the obstructions that separate a
collection of locally acceptable musical fragments from a
\emph{finished production}.

\paragraph{$H^0$: Global consistency.}
$H^0(\mathcal{C}_{\mathrm{DAW}}, \mathcal{F})$ is the space of
global sections---the data that is consistent across every open set.
Non-trivial $H^0$ elements include:
\begin{itemize}
  \item The master bus output is valid audio (finite amplitude,
        non-silent, within 0\,dBFS).
  \item All sample references in all instruments resolve to
        existing files.
  \item All plugin identifiers (\texttt{VST2}, \texttt{VST3},
        \texttt{AU}) resolve to installed binaries.
  \item BPM and time signature are consistent across the entire
        song (or the tempo automation is continuous).
\end{itemize}

\paragraph{$H^1$: Obstructions to gluing.}
$H^1 \neq 0$ is where the real problems live---where locally
acceptable sections fail to compose into a coherent whole.  Each
class of $H^1$ corresponds to a specific production problem:

\begin{enumerate}
  \item \textbf{Phase cancellation.}\;
        Two tracks---e.g.\ a layered snare top and snare
        bottom---that sound full in solo but partially cancel when
        summed to the mix bus.  This is a textbook $H^1$ class:
        $\mathcal{F}|_{U_{\mathrm{top}}}$ and
        $\mathcal{F}|_{U_{\mathrm{bot}}}$ are each satisfactory,
        but $\mathcal{F}|_{U_{\mathrm{top}} \cap
        U_{\mathrm{bot}}} = \mathcal{F}|_{\mathrm{sum}}$ has a
        non-trivial obstruction.
  \item \textbf{Frequency masking.}\;
        A kick drum and a bass synth both occupy the
        60--120\,Hz band.  Each sounds powerful in solo; together
        they produce mud.  The overlap
        $U_{\mathrm{kick}} \cap U_{\mathrm{bass}}$ in the
        spectral domain is the source of the obstruction.
  \item \textbf{Key clashes at pattern boundaries.}\;
        Pattern~$A$ is in C~major; pattern~$B$ modulates to
        F$\sharp$~major.  Each pattern is harmonically
        well-formed, but the transition---the overlap in the
        covering sieve---carries an $H^1$ class measuring the
        smoothness of the modulation.
  \item \textbf{Automation discontinuities.}\;
        A filter cutoff automation in pattern~$A$ ends at
        value~$x$; the same parameter in pattern~$B$ begins at
        value~$y \neq x$.  At the pattern boundary, the
        presheaf has a discontinuity---an $H^1$ obstruction in
        the automation sub-presheaf.
  \item \textbf{Sustained-note truncation.}\;
        A note sustains across a pattern break via Renoise's
        \texttt{OFF}/\texttt{---} convention, but the next
        pattern's first line contains a conflicting note-on in
        the same column.  The gluing fails.
  \item \textbf{Sidechain conflicts.}\;
        Two instruments sidechained to the same source
        (e.g.\ kick) with conflicting ducking envelopes: one
        demands fast release, the other demands slow release.
        The shared dependency creates an obstruction in the
        signal-flow presheaf.
  \item \textbf{Tempo automation breaks.}\;
        Pattern at 120~BPM jumps to 140~BPM without a
        ritardando or accelerando transition.  The tempo
        presheaf is discontinuous.
\end{enumerate}

\paragraph{$H^2$ and higher: Structural obstructions.}
Higher cohomology captures arrangement-level issues visible only in
full-song context:
\begin{itemize}
  \item The verse--chorus--verse--chorus form lacks a satisfying
        energy arc (no build, no release).
  \item Mix balance is locally fine section-by-section but the
        overall spectral center of gravity shifts
        unpredictably across the arrangement.
  \item Motivic development stalls---themes introduced in the
        intro never return or transform.
\end{itemize}

\begin{proposition}[Mayer--Vietoris for mixing]
\label{prop:mv-mixing}
Let $A$ and $B$ be two track groups (e.g.\ drums and bass).  Their
spectral/stereo overlap is $A \cap B$.  The Mayer--Vietoris
sequence
\[
  \cdots \to H^0(A \cup B) \to H^0(A) \oplus H^0(B)
  \to H^0(A \cap B) \xrightarrow{\;\partial\;}
  H^1(A \cup B) \to \cdots
\]
gives the connecting homomorphism $\partial$ that sends ``local
spectral compatibility'' to ``global mixing obstruction.''
Concretely: $A$ and $B$ each pass their solo mix checks
($H^0(A), H^0(B)$ trivial), but their shared frequency content
$H^0(A \cap B)$ maps via $\partial$ to a non-trivial $H^1(A \cup
B)$ class---frequency masking.  The standard mixing repair is to
apply complementary EQ curves to $A$ and $B$, carving space in
$A \cap B$ until $\partial = 0$.
\end{proposition}

% ────────────────────────────────────────────────────────────────
\subsection{The graded evaluation \texorpdfstring{$\GEval$}{Q} for
  music}
\label{ssec:musical-grading}

The graded evaluation for a DAW project is multi-stratal, with
each stratum measuring an independent axis of production quality.

\begin{definition}[Musical grade algebra]
\label{def:musical-grade}
Define the grade monoid
$G_{\mathrm{DAW}} = (G, \otG, \gone, \opG, \gzero)$ with strata:
\begin{center}
\renewcommand{\arraystretch}{1.2}
\begin{tabular}{@{}lll@{}}
\textbf{Stratum} & \textbf{What it measures} &
  \textbf{Computation} \\
\hline
$\mathsf{mix}$         & Mix quality
  & LUFS loudness, dynamic range, frequency balance, stereo width \\
$\mathsf{master}$      & Mastering quality
  & Peak level, RMS, crest factor, spectral tilt \\
$\mathsf{arrangement}$ & Structural quality
  & Section lengths, contrast ratios, tension/release arc \\
$\mathsf{harmony}$     & Harmonic quality
  & Voice leading, progression logic, modulation smoothness \\
$\mathsf{rhythm}$      & Rhythmic quality
  & Groove consistency, syncopation interest, quantize accuracy \\
$\mathsf{timbre}$      & Timbral quality
  & Spectral diversity, band allocation, stereo separation \\
$\mathsf{melody}$      & Melodic quality
  & Contour balance, interval distribution, motivic development \\
$\mathsf{production}$  & Production polish
  & Noise floor, artifact absence, latency compensation \\
\end{tabular}
\end{center}
The overall grade is the $\otG$-product across all strata:
\[
  g_{\mathrm{total}} = g_{\mathsf{mix}} \otG\, g_{\mathsf{master}}
    \otG\, g_{\mathsf{arrangement}} \otG\, g_{\mathsf{harmony}}
    \otG\, g_{\mathsf{rhythm}} \otG\, g_{\mathsf{timbre}}
    \otG\, g_{\mathsf{melody}} \otG\, g_{\mathsf{production}}.
\]
\end{definition}

The $\otG$-product enforces the \textbf{bottleneck principle}: a
track with brilliant melodies ($g_{\mathsf{melody}} \approx
\gone$) but a clipping, distorted mix ($g_{\mathsf{mix}} \approx
\gzero$) receives a low total grade, because $\gone \otG \gzero
= \gzero$.  This is the formal expression of the producer's maxim
that a great song with a bad mix never gets played.

\begin{example}[Loudness war as grade collapse]
The ``loudness war'' is the phenomenon of over-compressing masters
to maximize perceived loudness at the cost of dynamic range.  In
our grading, aggressive brick-wall limiting pushes
$g_{\mathsf{master}}$ toward $\gone$ (the LUFS target is met) but
simultaneously drives $g_{\mathsf{mix}} \to \gzero$ (dynamic range
collapses, transients are destroyed).  The $\otG$-product correctly
identifies this as a degradation: the master stratum cannot
compensate for a destroyed mix stratum.
\end{example}

% ────────────────────────────────────────────────────────────────
\subsection{Descent as the production workflow}
\label{ssec:musical-descent}

The actual DAW production workflow is \emph{descent}: a sequence of
refinements that reduce $H^1$ obstructions and increase $\GEval$
grades.

\begin{enumerate}
  \item \textbf{Sketch (initial candidate).}\;
        Rough patterns with placeholder sounds---a bare melodic
        idea in a default sine wave, a four-on-the-floor kick from
        the factory library.  This candidate has many $H^1$
        obstructions: no frequency separation, no arrangement
        logic, placeholder timbres.
  \item \textbf{Arrangement (first descent step).}\;
        Organize patterns into a song structure---intro, build,
        drop, break, drop, outro.  Assign section types from the
        musical type $\tau$.  This resolves structural $H^1$
        classes (the piece now has a coherent macro-form) and
        increases $g_{\mathsf{arrangement}}$.
  \item \textbf{Sound design (second step).}\;
        Replace placeholder instruments with designed sounds:
        layer samples in Renoise's instrument editor, sculpt
        wavetables, program modulation sets and macros.  This
        resolves timbral $H^1$ (the sounds now occupy distinct
        spectral regions) and increases $g_{\mathsf{timbre}}$.
  \item \textbf{Mixing (third step).}\;
        Apply EQ, compression, saturation, and spatial
        processing.  Configure send tracks for reverb and delay.
        Set up sidechain routing.  Automate levels across
        sections.  This resolves frequency-masking, phase, and
        stereo $H^1$ classes, and increases $g_{\mathsf{mix}}$.
  \item \textbf{Mastering (final step).}\;
        Master bus processing: multi-band compression, stereo
        imaging, limiting to target LUFS.  Resolves loudness and
        dynamic-range $H^1$, and pushes $g_{\mathsf{master}}$
        toward $\gone$.
\end{enumerate}

Each step is one iteration of the descent loop from
Definition~\ref{def:musical-site}: produce a candidate section,
compute $\CSpec$ (detect obstructions), compute $\GEval$ (grade
quality), repair the lowest-grade stratum, and repeat.

\begin{remark}[Renoise-specific descent moves]
Renoise's pattern effect commands are \emph{local repair
operations}:
\begin{itemize}
  \item \texttt{0Exx} (retrigger) repairs rhythmic sparseness by
        subdividing a single note event.
  \item \texttt{0Cxx} (set volume) repairs dynamic-range
        obstructions within a pattern.
  \item \texttt{0Gxx} (glide-to-note) repairs melodic
        discontinuities by interpolating between pitches.
  \item DSP chain reordering---EQ before vs.\ after
        compression---changes the $H^1$ structure: pre-EQ
        tightens the compressor's response (repairs spectral
        obstructions first), while post-EQ allows broader
        dynamics but may reintroduce masking.
  \item Instrument macro mappings automate multi-parameter descent:
        a single macro knob can simultaneously adjust filter
        cutoff, resonance, and envelope decay, maintaining
        timbral consistency as one control sweeps a family of
        coherent states.
\end{itemize}
\end{remark}

% ────────────────────────────────────────────────────────────────
\subsection{Music generation via \texorpdfstring{$\FMLM$}{FMLM}}
\label{ssec:musical-fmlm}

JuGeo can \emph{generate} music, not just verify it.  The
$\FMLM$ loop (Fill--Measure--Locate--Mend) applies directly.

\begin{example}[Techno track generation]
\label{ex:techno-gen}
Start with a type:
\[
  \tau = \mathsf{Song}\!\left(
    \begin{aligned}
      &\mathit{genre} = \mathsf{techno},\;
       \mathit{bpm} = 128,\;
       \mathit{key} = \mathrm{A}\,\text{minor}, \\
      &\mathit{structure} = [\mathsf{intro}, \mathsf{build},
        \mathsf{drop}, \mathsf{break}, \mathsf{drop},
        \mathsf{outro}]
    \end{aligned}
  \right).
\]
The $\FMLM$ loop proceeds:
\begin{enumerate}
  \item \textbf{Fill (iteration~1): drums.}\;
        For each section, fill the drum track patterns---kick,
        hats, clap, ride---respecting the rhythmic constraints
        of $\tau$.  After this step, the project has rhythmic
        content but massive $H^1$ (no harmonic content, no
        spectral balance).
  \item \textbf{Measure:}\;
        compute
        $g_{\mathsf{rhythm}} \approx 0.7$,
        $g_{\mathsf{harmony}} = \gzero$,
        $g_{\mathsf{timbre}} \approx 0.2$.
  \item \textbf{Fill (iteration~2): bass line.}\;
        Add a bass track constrained by the key
        ($\mathrm{A}\,\text{minor}$) and the groove (locked to
        kick pattern).  Frequency band: 30--200\,Hz.
  \item \textbf{Measure:}\;
        $g_{\mathsf{harmony}} \approx 0.4$ (bass provides root
        motion), $H^1_{\mathrm{mask}}$ appears between kick and
        bass in the 50--80\,Hz band.
  \item \textbf{Locate + Mend:}\;
        apply sidechain compression from kick to bass
        (resolve masking $H^1$), high-pass the bass at 35\,Hz.
  \item \textbf{Fill (iteration~3): lead/pad.}\;
        Add a pad track (200\,Hz--2\,kHz) and a lead
        (1--8\,kHz), constrained by the harmonic type and the
        remaining spectral space.
  \item \textbf{Fill (iteration~4): FX and transitions.}\;
        Add risers, sweeps, impacts at section boundaries.
        Automate filter cutoffs for builds.  Resolve
        arrangement-level $H^1$ (section transitions are now
        smooth).
  \item \textbf{Final Measure:}\;
        all strata near $\gone$; $H^1 \approx 0$; the project
        is a well-typed, cohomologically trivial,
        high-grade \texttt{.xrns} file.
\end{enumerate}
This is \emph{exactly} how an experienced tracker musician
works---JuGeo formalizes the workflow.
\end{example}

% ────────────────────────────────────────────────────────────────
\subsection{The isomorphism: music \texorpdfstring{$\cong$}{≅}
  software}
\label{ssec:musical-iso}

The structural parallel between a DAW project and a software
project is not a loose analogy but a genuine categorical
equivalence mediated by the 4-tuple.

\begin{proposition}[DAW--software equivalence]
\label{prop:daw-software}
There exists a structure-preserving functor
$\Phi \colon \mathcal{C}_{\mathrm{DAW}} \to
\mathcal{C}_{\mathrm{SW}}$ that maps:
\begin{center}
\renewcommand{\arraystretch}{1.2}
\begin{tabular}{@{}ll@{}}
\textbf{Music (DAW project)} & \textbf{Software project} \\
\hline
Track              & Module / thread \\
Pattern            & Function / procedure \\
Note / effect cmd  & Instruction / statement \\
Signal flow        & Data flow \\
Mixer routing      & Dependency graph \\
Plugin (VST/AU)    & Library (npm, pip, crate) \\
Frequency masking  & Resource contention \\
Phase cancellation & Race condition \\
Clipping (0\,dBFS) & Buffer overflow \\
Sidechain routing  & Semaphore / lock \\
Mastering          & Deployment / release \\
Mix balance        & Load balancing \\
Automation curve   & Configuration schedule \\
Pattern matrix     & Module dependency matrix \\
\end{tabular}
\end{center}
Under $\Phi$, the 4-tuples are compatible:
$\Phi_*(\STop_{\mathrm{DAW}}) = \STop_{\mathrm{SW}}$,
$\Phi_*(\PStr_{\mathrm{DAW}}) = \PStr_{\mathrm{SW}}$, and
similarly for $\CSpec$ and $\GEval$ (with appropriate stratum
relabeling).
\end{proposition}

\begin{remark}
The equivalence is productive in both directions.  Software
engineering concepts suggest musical techniques: \emph{unit
testing} corresponds to \emph{solo monitoring} (check each track
in isolation), \emph{integration testing} to \emph{group bus
monitoring} (check track groups together), and \emph{end-to-end
testing} to \emph{full-mix playback}.  Conversely, musical
concepts suggest software techniques: \emph{sidechain
compression}---dynamically ducking one signal in response to
another---is a pattern for adaptive resource throttling, and
\emph{send/return routing}---sharing a single reverb instance
across many tracks---is a service mesh pattern.
\end{remark}

\begin{remark}
Algorithmic composition has a long history, from Hiller and
Isaacson's \emph{Illiac Suite}~(1957) to modern neural audio
synthesis.  What JuGeo adds is the cohomological dimension: not
merely ``does each pattern sound acceptable in isolation?''
($H^0$) but ``do the patterns compose into a coherent
production?'' ($H^1$).  This is precisely the difference between
a collection of loops and a \emph{finished track}---and it is the
same difference that separates a collection of working functions
from a \emph{shipped program}.
\end{remark}

%% ─── 30.5 ──────────────────────────────────────────────────────
\section{Ethical Judgment Geometry}
\label{sec:ethical}

Ethical reasoning involves multiple, potentially conflicting moral
frameworks---a natural setting for graded, multi-theory judgment.

\begin{itemize}
  \item \textbf{Moral frameworks as theory slots.}\;
        Utilitarianism, deontology, virtue ethics, care ethics,
        and other frameworks each become a theory slot in the CTT.
        Each provides a grade for a given action or policy.
  \item \textbf{Ethical dilemmas as $H^1 \neq 0$.}\;
        An ethical dilemma is precisely a situation where different
        moral frameworks assign contradictory verdicts: the local
        sections (framework-specific judgments) fail to glue.
        The resulting $H^1$ class identifies the \emph{source} of
        the disagreement.
  \item \textbf{Resolution via graded evaluation.}\;
        Rather than forcing a binary choice, $\GEval$ assigns a
        grade that reflects the degree of moral consensus or
        conflict.  A policy with $g \approx \gone$ enjoys broad
        agreement; a policy with $g \ll \gone$ is deeply contested.
  \item \textbf{Pluralism via the open world property.}\;
        New ethical frameworks can be added to the theory ensemble
        without modifying existing ones.  This is genuine moral
        pluralism: no framework is privileged, and the graded
        evaluation tracks how well any action satisfies the full
        ensemble.
\end{itemize}

%% ─── 30.6 ──────────────────────────────────────────────────────
\section{JuGeo for education}
\label{sec:education}

Educational assessment is a domain where graded judgment is
the norm, not the exception.

\begin{itemize}
  \item \textbf{Student work as evidence.}\;
        A student's submission $E$ is evidence for a rubric
        specification~$A$.  The judgment
        $\vdash_g E : A$ says ``$E$ satisfies $A$ to grade~$g$\!.''
  \item \textbf{Partial credit.}\;
        The $\Grade$ semiring provides a principled theory of partial
        credit.  Each rubric criterion is a theory slot with its own
        grade; the total grade is $\otG$-aggregated (all criteria
        matter) or $\opG$-aggregated (best-of, for alternative
        approaches).
  \item \textbf{Feedback as obstruction identification.}\;
        When $g < \gone$, the obstructions identify which
        rubric criteria were not fully met.  This is structured
        feedback: not just ``your grade is 7/10'' but ``the
        logical structure criterion is at $-0.2$ because of a gap
        in the proof of Lemma~3.''
  \item \textbf{Automated grading via descent.}\;
        The grading process itself can be automated:
        parse the submission, compute theory-slot grades, aggregate,
        and report obstructions.  This is structurally identical to
        program verification---the student's work is a ``program''
        and the rubric is its ``specification.''
\end{itemize}

%% ─── 30.7 ──────────────────────────────────────────────────────
\section{Self-improving JuGeo}
\label{sec:self-improving}

In the spirit of G\"odelian self-reference, JuGeo can be applied
to its own development.

\begin{construction}[JuGeo Bootstrap]
\label{con:bootstrap}
Let $\STop_{\mathrm{JG}}$ be the site of JuGeo's own codebase
(the modules in \texttt{src/jugeo/}).  Define:
\begin{itemize}
  \item $\PStr_{\mathrm{JG}}$: JuGeo's specification (the
        mathematical definitions of this monograph) together with
        the implementation as evidence.
  \item $\CSpec_{\mathrm{JG}}$: the \v{C}ech cohomology of the
        implementation relative to the specification.
        $H^1 \neq 0$ identifies modules whose implementation
        deviates from the formal theory.
  \item $\GEval_{\mathrm{JG}}$: the graded evaluation of JuGeo
        by JuGeo---using the very descent algorithm being verified
        to check that the descent algorithm is correctly
        implemented.
\end{itemize}
\end{construction}

This bootstrap is not circular: JuGeo version $n$ verifies JuGeo
version $n+1$, and each version's correctness is established
relative to the previous one.  The process is analogous to
compiler bootstrapping, where compiler version $n$ compiles
version~$n+1$.

%% ─── 30.8 ──────────────────────────────────────────────────────
\section{Open mathematical problems}
\label{sec:open-problems}

We close with a list of open problems that the framework raises.

\begin{enumerate}
  \item \textbf{Decidability.}\;
        Is the full 4-tuple verification problem ($\check{H}^1 = 0$
        and $\GEval \geq g_{\min}$) decidable for any interesting
        fragment of the site category?  For finite sites with
        computable grade functions the answer is trivially yes; for
        sites arising from Turing-complete programming languages
        it is trivially no.  The interesting question is: what is
        the largest class of sites for which the problem remains
        decidable, and does this class cover practically relevant
        programs?

  \item \textbf{Homotopy type of the judgment fibration.}\;
        The judgment fibration $\pi\colon \mathcal{E} \to \STop$
        is a Grothendieck fibration with rich structure.  What are
        its homotopy invariants?  In particular, is the classifying
        space $B\!\operatorname{Aut}(\mathcal{E})$ related to known
        classifying spaces in algebraic topology?

  \item \textbf{Spectral sequences for NLP.}\;
        The \v{C}ech-to-derived spectral sequence converging to
        $H^*(\STop, \PStr)$ provides a computational tool that, to
        our knowledge, has never been applied to NLP\@.  Can the
        differentials $d_r\colon E_r^{p,q} \to E_r^{p+r,q-r+1}$
        be given linguistic interpretations?  Could this yield new
        parsing or generation algorithms that operate stratum-by-stratum?

  \item \textbf{Continuous JuGeo.}\;
        Replace the discrete site with a smooth manifold (or a
        diffeological space).  Sections become smooth functions,
        the \v{C}ech complex becomes the de Rham complex, and
        obstructions become differential forms.  Does this
        ``continuous JuGeo'' yield a useful framework for
        cyber-physical systems, control theory, or signal processing?

  \item \textbf{Categorification of the grade.}\;
        The $\Grade$ semiring is a 0-categorical structure.  Can it
        be categorified---replacing grades by grading categories,
        and the semiring operations by functors---to capture
        ``grades of grades''?  This would give a natural notion of
        \emph{confidence in the grade itself}, addressing the
        meta-question ``how reliable is our quality assessment?''

  \item \textbf{Effective descent bounds.}\;
        Proposition~\ref{prop:llm-descent-mono} gives a worst-case
        bound on descent iterations.  What is the \emph{expected}
        number of iterations under realistic distributions of LLM
        errors?  Is there a phase-transition phenomenon: a critical
        error rate above which descent fails to converge?

  \item \textbf{JuGeo and dependent type theory.}\;
        Part~III introduced graded dependent types.  To what extent
        can JuGeo be \emph{internalized} in a dependent type theory
        (e.g., Agda or Lean~4)?  A full internalization would yield
        machine-checked proofs of JuGeo's own metatheory.
\end{enumerate}

\bigskip

These problems span category theory, homotopy theory, computability,
and applied mathematics.  We hope that the framework developed in
this monograph---simple enough to be stated in a single 4-tuple,
rich enough to organize program verification, natural language
generation, and beyond---will inspire progress on all of them.
